%%%%%%%%%%%%%%%%%%%%%%%%%%%%%%%%%%%%%%%
\chapter{My approach}
\label{chap:approach}
%%%%%%%%%%%%%%%%%%%%%%%%%%%%%%%%%%%%%%%


Après l'état de l'art, les parties varient en fonction de votre
recherche. Vous devez adapter vos titres : Problématique,
Développement, Méthodologie, Résultats\ldots


Ces parties vont servir à décrire précisément le développement de
votre recherche : ce que vous avez fait, comment vous vous y êtes
pris, pourquoi, les résultats qui en découlent, les analyses que vous
en faites. C'est dans cette partie que vous allez plonger le lecteur
dans votre démarche scientifique. Cette partie peut être variable d'un
mémoire à l'autre, en fonction du sujet étudié, de la recherche menée
ou développée, des attentes de votre promoteur. Quoi qu'il en soit,
généralement, il est attendu que l'on retrouve les éléments suivants
dans cette partie :

\begin{itemize}
\item[\textbullet] une description de votre problématique/question de
  recherche ;
\item[\textbullet] une description et une justification de votre
  méthodologie (participants ou public-cible visé, étapes de
  conception et/ou de testing, méthode de collecte de données, de
  description des analyses, …) ;
\item[\textbullet] une description de vos résultats : il s'agit bien
  d'une description des données récoltées ou de l'outil, de
  l'application, … Vous êtes dans une écriture objective, il n'y a
  aucune interprétation personnelle.
\end{itemize}


\paragraph{Langue.} Ce template est configuré par défaut pour la
langue française. Si votre mémoire est en anglais, songez à éditer la
ligne avec la commande \texttt{\textbackslash selectlanguage}.

\paragraph{Sauts de ligne.}
\'Evitez d'utiliser des sauts de lignes dans un paragraphe ou
d'utiliser un paragraphe par phrase. Un paragraphe est un ensemble de
phrases réunies par une même unité de pensée.\\ Et on constate ici
qu'un saut de ligne est bien disgracieux\dots.

\paragraph{Césures.} \LaTeX\ peut parfois éprouver des difficultés pour opérer une césure
de mots longs. On peut alors l'aider en plaçant des indications de
césures comme dans \og le mot le plus plus long est
an\-ti\-cons\-ti\-tu\-tion\-nel\-le\-ment\fg{} qui a été édité comme :
\texttt{an\textbackslash-ti\textbackslash-cons\textbackslash-ti\textbackslash-tu\textbackslash-tion\textbackslash-nel\textbackslash-le\textbackslash-ment}.

\paragraph{Guillemets} Le package \texttt{babel} pour la langue
française fournit des commandes \verb+\og+ et \verb+\fg{}+ pour
produire des guillemets conformes aux règles de typographie française~:
\og Entre guillemets, \fg{} c'est plus joli que ceci "assurément". Les
accolades fermantes sont nécessaires après \verb+\fg+ afin d'empêcher
la commande d'absorber le caractère suivant comme un argument.

Le package \texttt{csquotes} utilise une commande \verb+\enquote+ avec
argument qui produit des guillemets conformes à la langue
utilisée. Ainsi, ce code \LaTeX :
\begin{verbatim}
\setquotestyle{english}
  This is an \enquote{English} paragraph.
\setquotestyle{french}
  Ceci est un paragraphe en \enquote{Français}.    
\end{verbatim}

produira:

\begin{quote}
  This is a quote.
\end{quote}

Il est toujours possible d'utiliser pour la langue anglaise les
caractères \verb+``+ (accent grave) et \verb+''+ (quote) qui
produisent ``entre guillemets''.

Dans tous les cas, l'usage des guillements simples \og \texttt{"}
\fg{} est à proscrire.


\section{Les figures et tables}
\label{sec:figandtab}

Une figure est un corps flottant dans votre document. \LaTeX\ se
charge par défaut de le placer au meilleur endroit afin de privilégier
l'esthétique de vos pages et éviter des blancs disgracieux. Il est
souhaitable de guider \LaTeX\ pour qu'il place vos figures de
préférence en haut ou sur une page séparée (paramètre
\texttt{[tp]}). Cela produit les meilleurs résultats. Il est également
possible de forcer le positionnement (paramètre \texttt{[h]}), mais
cette politique doit être utilisée avec parcimonie, pour la raison
énoncée plus haut.

Toutes les figures (et le tables) doivent être référencées
explicitement dans le texte avec la commande
\verb+\ref{fig:mafigure}+. Vous pouvez utiliser la commande
\verb+\pageref{fig:mafigure}+ pour indiquer le numéro de page. Cela
peut être utile si votre figure est fort éloignée, dans les annexes
par exemple.  Pour faire référence à une figure dans le texte, il
suffit de l'indiquer entre parenthèses dans le corps de texte comme
ceci (Figure \ref{fig:mafigure}, page~\pageref{fig:mafigure}, dans la
section~\ref{sec:figandtab}) ou de manière explicite: \textit{la
  figure~\ref{fig:mafigure} illustre bla bla}.

\begin{figure}[ht]
	\centering
        %% vous pouvez utiliser width=\textwidth pour ajuster la figure à la largeur de la feuille, ou utiliser un facteur multiplicatif pour une largeur proportionnelle, comme width=0.33\textwidth
        \includegraphics[height=5cm]{img/figure1.jpg}
        \caption{Titre de la figure}
        \label{fig:mafigure} 
\end{figure}

\LaTeX\ fait une distinction entre les Figures et les Tables. Chacune
dispose par défaut d'une numérotation spécifique et d'une mise en page
différente. Il est préférable de délaisser les environnements Table
afin de privilégier une certaine homogénéité.

 \section{Algorithme}


 Si vous devez produire du code, vous pouvez utiliser le package
 \texttt{listings}. Par défaut, ce package reconnaît de nombreux
 langages de programmation.

\begin{lstlisting}[language=Java,caption=Classe HelloWorld en Java]
public class HelloWorld {
    public static void main(String[] args) {
        // Affiche un message
        System.out.println("Bonjour, monde !");
    }
}
\end{lstlisting}

Il est également possible (voire préférable) d'externaliser le code
dans un fichier spécifique et de l'insérer avec la commande
\texttt{lstinputlisting}.

\begin{verbatim}
\lstinputlisting[caption=HelloWorld.java]{HelloWorld.java}
\end{verbatim}

Pour des langages plus ésotériques ou des besoins plus spécifiques, le
package \texttt{minted} offre plus de possibilités, mais il introduit
une dépendance supplémentaire en soutraitant la compilation à
Python. Il est donc plus lent et son installation est plus
compliquée. Le code source de votre document est moins portable
également.
 
 \section{Pseudo-Code}

 Si vous devez produire du pseudo-code, l'environnement
 \texttt{algorithm2e} est sans doute le plus adapté.

 \begin{algorithm}[H]
\DontPrintSemicolon
\KwData{Un tableau $A$ de $n$ entiers}
\KwResult{Le plus grand élément de $A$}
$max \leftarrow A[1]$\;
\For{$i \leftarrow 2$ \KwTo $n$}{
    \If{$A[i] > max$}{
        $max \leftarrow A[i]$\;
    }
}
\Return $max$\;
\caption{Recherche du maximum dans un tableau}
\end{algorithm}